\section{Rare events and the inverse-square boundary}\label{sec:rare}
\subsection{The rare-event experiment}\label{sec:rare-model}
The preceding minimax law takes a worst case over every sign function.
For rare events, fix $0<\rho\le1/4$ and restrict this class to
\[
 \mathcal Q_\rho(\mu)=
 \left\{q:(0,1/2]\to[0,1]\text{ measurable}:\int q\,d\mu\le\rho\right\}.
\]
We write $\mathcal Q_\rho$ when $\mu$ is fixed.
The unknown target is still $\theta(q)=\int q\,d\mu$; $\rho$ is a known
upper bound. For every $B>0$ define
\[
 R_{\mu,\rho}(B)=
 \inf_{\mathcal A:\,\sup_{q\in\mathcal Q_\rho(\mu)}\E_qT_\mathcal A\le B}
 \sup_{q\in\mathcal Q_\rho(\mu)}
       \bigl(\E_q(\widehat\theta_\mathcal A-\theta(q))^2\bigr)^{1/2}.
\]
The observation model remains the ordinary iid fresh-coin experiment:
each fresh unit samples a hidden magnitude $D\sim\mu$ and a hidden sign
with conditional positive probability $q(D)$; its observations are iid
Bernoulli with mean $1/2\pm D$ according to the sign. Different units
have independent latent states and observation tapes, and a unit keeps
its latent state on every revisit. Individual magnitudes and signs are
hidden; all creation, revisiting, interleaving, and stopping decisions
use only an independent seed and the observed flips. Every requested
observation costs one, including the first on a unit; units created but
never observed have no statistical effect. Stopping may be unbounded,
and finite expected work gives almost-sure termination under every
admissible population.

The benchmark may know $\mu,\rho,B$. In particular, the global constraint
does not impose $q(d)\le\rho$ pointwise, nor fix either conditional
magnitude law given the sign. Estimates may be clipped to $[0,\rho]$
without increasing risk. Returning $\rho/2$ costs nothing and has
RMS error at most $\rho/2$; randomization permits meaningful procedures
even for $B<1$. The power-law and critical results below are stated for
$B\ge1$, as in their sources; Sections~\ref{sec:entropy-tail}
and~\ref{fg:sec:model} cover every $B>0$.
\Cref{sec:entropy-tail,sec:certificates} and their appendices write the
sign as $S\in\{0,1\}$, with observation mean $1/2+(2S-1)D$;
\cref{fg:sec:model} and \cref{app:finite-direct} write $S\in\{-1,+1\}$,
with mean $1/2+SD$.

Two different boundaries occur in this manuscript. The unrestricted-sign
inverse-moment criterion changes at power exponent $\alpha=1$; the
rarity-dependent quadratic phenomenon below occurs at $\alpha=2$.

\subsection{A variational lower modulus and its exact optimizer}
For $h>0$ set
\begin{equation}\label{eq:v5-psi}
 \Psi_{\mu,\rho}(h)=
 \sup_{\substack{v\text{ measurable},\ 0\le v\le1\\\int v\,d\mu\le\rho}}
       \int\min\{v(d),h\sqrt{v(d)}/d\}\,\mu(dd).
\end{equation}
There is a universal $c>0$ such that
$R_{\mu,\rho}(B)\ge c\Psi_{\mu,\rho}(B^{-1/2})$.
The corresponding confidence lower bound replaces $B^{-1/2}$ by
$\sqrt{\log(1/\delta)/B}$. The proof in \cref{app:localization}
explicitly imports the single-coin rare-prior Hellinger bound of
Lee--Valiant~\cite[Corollary 19]{leevaliant}; its bounded-ratio hybrid
argument is written out. These statistical lower bounds are paper
proofs with that published external input. The general-law lower proof
in \cref{sec:matching-lower} does not use that input: it derives the
null-work consequence of $\Psi$ directly in
\cref{sec:transfer-foundations}.

The variational optimization itself is fully explicit. There is a unique
$t\in(0,1/2)$ with
\[
 \E_\mu\min\{1,t^2/D^2\}=\rho.
\]
For every $h>0$, $v_*(d)=\min\{1,t^2/d^2\}$ is an optimizer. For $h\ge t$
the value is $\rho$; for $h<t$ it is
\begin{equation}\label{eq:v5-psi-form}
 F_\mu(h)+h\E[D^{-1};\,h<D\le t]+ht\E[D^{-2};\,D>t].
\end{equation}
Atoms and singular laws are allowed, with no inverse-moment assumption.
The Lean theorem
\leanref{Tier3/Optimizer.lean}{arbitrary\_law\_optimizer} proves this
functional statement. It does not supply an attaining statistical
algorithm. In particular,
\[
 \Psi_{\delta_a,\rho}(h)=\min\{\rho,h\sqrt\rho/a\},\qquad
 \rho\omega_\mu(h/\sqrt\rho)\le\Psi_{\mu,\rho}(h)\le\omega_\mu(h).
\]

\subsection{The power-margin regimes}
Let $\mu_{\alpha,a}$ have distribution function $(d/a)^\alpha$ on
$(0,a]$ and put $x=(a\sqrt B)^{-1}$. The proved RMS orders are:
\begin{center}
\renewcommand{\arraystretch}{1.4}
\begin{tabular}{@{}ll@{}}
\toprule
Magnitude law & RMS order, with saturation included\\
\midrule
$0<\alpha<1$ & $\min\{\rho,x^\alpha\}$\\
$\alpha=1$ & $\min\{\rho,x[1+\log_+(\rho/x)]\}$\\
$1<\alpha<2$ & $\min\{\rho,\rho^{\,1-1/\alpha}x\}$\\
$\alpha=2$, fixed $a=1/4$
 & $\min\{\rho,\sqrt{\rho L(1+\log L)/B}\}$\\
$\alpha>2$ & $\min\{\rho,\sqrt\rho\,x\}$\\
\bottomrule
\end{tabular}
\end{center}
Here $L=1+\log(1/\rho)$. For the rows below two, $0<a\le1/4$;
for the row above two, $0<a\le1/2$. Constants may depend on the fixed
exponent. The quadratic row concerns one normalized law, with universal
constants; it does not claim a theorem uniform in $\alpha$ near two.
All rows hold for $B\ge1$ and $0<\rho\le1/4$ in their stated domains.

For $0<\alpha<2$, \cref{app:localization} gives full proofs, including
confidence radii after replacing $x$ by
$a^{-1}\sqrt{\log(1/\delta)/B}$. The lower bound restricts the balanced
hard alternatives to a region of magnitude mass $\rho$ and does not
use the imported rare-prior theorem. The upper bound starts the
randomized telescope at a large enough cap that its base false-positive
probability is at most $\rho$, then uses the same two-success
calibration as \cref{sec:minimax-proof}. It knows $(\alpha,a,\rho)$.
It also holds under the envelope
$F_\mu(d)\le\min\{(d/a)^\alpha,1\}$.

For $\alpha>2$, \cref{app:above-two} supplies an asymmetric sequential
classifier: its negative-coin error is at most $\rho\varepsilon$, while
its positive-coin error is at most $\varepsilon$. A randomized telescope
then yields a one-coin unbiased output with separate bounds
$\sup_q\E Y^2\le(153/56)\rho$ and $\sup_q\E T_Y\le w_\alpha/a^2$,
where $w_\alpha<\infty$ is explicit there. Independent Poisson replication
calibrates every real budget, including expected replication counts
below one. The upper bound is uniform over the same envelope class and
knows $(\alpha,a,\rho,B)$.
The matching exact-power lower bound uses the published fixed-gap
Lee--Valiant theorem through a simulation that keeps a coin's privately
sampled magnitude fixed on every revisit. Constants diverge as
$\alpha\downarrow2$. No sharp confidence characterization above or at
two is asserted in this manuscript.

Both bounds of the row above two are also Lean-checked in the observable
model, by a different construction:
\abovetwoleanref{above\_two\_theory} combines a single-coin family that
depends only on $(a,\rho)$, is unbiased on every power-margin law, and
has second moment at most $28\rho$, with fixed-count averaging and a
lower bound with coefficient $1/256$ for every law supported below $a$.
The formal lower proof does not use the Lee--Valiant input. See
\cref{sec:verification} for the exact scope.

\begin{remark}[Agreement with the closed quantile form]\label{rem:power-quantile}
Each row also follows, with constants depending on $\alpha$, from
Corollary~\ref{fg:cor:general-laws}, which states
$R_{\mu,\rho}(B)\asymp\Q(B)$ for every law. Write $h=B^{-1/2}=ax$.
The quantile point of $\mu_{\alpha,a}$ is $t=a\rho^{1/\alpha}$. If
$h\ge t$, then $\Q=\rho$ and every row is at least a constant multiple
of $\rho$. If $h<t$, substituting $d=as$ gives
\[
 T_h=x^\alpha+\alpha x\int_x^{\rho^{1/\alpha}}s^{\alpha-2}\,ds,
\]
which is comparable to $x^\alpha$ for $\alpha<1$, equals
$x[1+\log(\rho/x)]$ for $\alpha=1$, and is at most a constant multiple of
$x\rho^{1-1/\alpha}$ for $\alpha>1$. The dyadic inverse-square costs of
$\kappa=\mu|_{[t,a]}$ decrease geometrically away from $t$ when
$\alpha<2$ and away from $a$ when $\alpha>2$, so their entropy is bounded
and $\Gamma_\kappa$ is comparable to $a^{-2}\rho^{1-2/\alpha}$ and $a^{-2}$,
respectively. At $\alpha=2$ the order-$L$ bands meeting $[t,a]$ carry
comparable costs, so $\Gamma_\kappa\asymp a^{-2}L(1+\log L)$. Hence
$h\sqrt{\rho\Gamma_\kappa}$ is comparable to $x\rho^{1-1/\alpha}$,
$x\sqrt{\rho L(1+\log L)}$ or $x\sqrt\rho$. Comparing the two terms of $\Q$
recovers each row; below exponent one, $x\rho^{1-1/\alpha}\le x^\alpha$
because $x<\rho^{1/\alpha}$. This is a consistency check between
derivations, not an additional verification: the corollary rests on
the paper proof of Theorem~\ref{thm:entropy-tail-law}.
\end{remark}

\subsection{The full critical theorem and the correction to \texorpdfstring{$\Psi$}{Psi}}
\begin{theorem}[Critical quadratic law]\label{thm:v5-critical}
Fix $\mu(dd)=32d\,dd$ on $(0,1/4]$. There are universal
$0<c<C<\infty$ such that for all $0<\rho\le1/4$ and $B\ge1$,
\[
 c\min\{\rho,\sqrt{\rho L H/B}\}
 \le R_{\mu,\rho}(B)
 \le C\min\{\rho,\sqrt{\rho L H/B}\},\qquad
 L=1+\log(1/\rho),\quad H=1+\log L.
\]
The lower bound covers every observable randomized population algorithm,
including nonlinear outputs, unbounded stopping, and arbitrary creation
and interleaving. The upper construction may know $\mu,\rho,B$.
\end{theorem}
\noindent The complete paper argument is in \cref{app:critical}.
Both bounds are Lean-checked in the original observable experiment:
\par\noindent
\leanref{CriticalEntropy/CriticalMinimax.lean}{critical\_minimax\_lower\_bound}
and \leanref{CriticalUpper/Full.lean}{critical\_minimax\_upper\_bound}.
The formal upper uses $C=100000$ and an alternative observable construction;
see \cref{sec:verification} for its exact scope.

The proof has three substantive transitions. First, sequential testing
at many possible gaps forces a scale-entropy cost. Conditional moments,
activation normalization, and repeated runs on unused flips of the
\emph{same} hidden coin transfer that fact to a small-bias one-coin
estimator: its variance supremum times its expected work under the
all-negative population is at least $c_0\rho LH$.
Second, a finite-experiment RMS comparison and compact convex minimax
argument obtain \emph{one} population pair for every capped single-coin
policy. A policy-dependent hard pair would not suffice.
Third, a finite-prefix midpoint decomposition transfers this pair to
interacting population algorithms, charging each flip once. The same
pair works at every cap; finite expected work under both endpoints
justifies passage to terminal outputs. The longer proof, including the
upper construction, is in \cref{app:critical}.

\paragraph{Relation to gap entropy.}
The sign-testing mechanism follows the line from Farrell~\cite{farrell}
and Chen--Li~\cite{chenli} to the distributional gap-entropy bound of
Chen--Li--Qiao~\cite[Lemma 4.1 and Remark D.4]{clq}, which includes a
Bernoulli extension. The finite-gap proof in \cref{app:critical} supplies
the required one-sided cost bound at fixed error without fair-coin
termination. Agrawal--Ram--Ramdas~\cite{agrawal} give general
composite-test bounds in the Farrell regime. The recent
Yao--Zhao--Zhou result~\cite{yao} resolves the gap-entropy conjecture for
Gaussian best-arm identification and reports Lean proofs; its target and
observation model differ from the population problem here.

At the quadratic law, annular mass of order $16^{-j}$ balances testing
cost of order $16^j$. Across order $L$ usable scales, comparable
normalized weights have entropy of order $\log L$. This explains $H$;
it does not supply a general-law theorem; \cref{sec:entropy-tail} obtains
one by a different argument. The contribution being assessed
is the null-cost obstruction and its transfer through one common hard
pair to every adaptive population algorithm. No novelty is claimed for
the underlying sign-entropy mechanism.

For this fixed law, the optimizer identity is
\[
 \rho=(t/a)^2[1+2\log(a/t)],\qquad
 \Psi_{\mu,\rho}(h)=
 \begin{cases}
 [2ht(1+\log(a/t))-h^2]/a^2,&0<h<t,\\
 \rho,&h\ge t,
 \end{cases}
 \quad a=1/4.
\]
The exact formula is proved by
\leanref{CriticalPsi/Full.lean}{critical\_psi\_formula} for the actual supremum
over all measurable feasible allocations. Moreover,
\leanref{CriticalPsi/Full.lean}{critical\_psi\_rate} proves, for $h>0$,
\[
 \min\{\rho,h\sqrt{\rho L}\}\le\Psi_{\mu,\rho}(h)
 \le16\min\{\rho,h\sqrt{\rho L}\}.
\]
Thus $\Psi_{\mu,\rho}(B^{-1/2})\asymp\min\{\rho,\sqrt{\rho L/B}\}$.
For the paper's explicit sequence, taking $\rho_n=\exp(-\exp n)$ and $B_n=\rho_n^{-2}$ gives
\[
 \frac{R_{\mu,\rho_n}(B_n)}
      {\Psi_{\mu,\rho_n}(B_n^{-1/2})}
 \ge c'\sqrt{1+\log(1+\log(1/\rho_n))}\longrightarrow\infty.
\]
Thus the proposed universal matching-$\Psi$ upper bound is false
according to \cref{thm:v5-critical}, even for one known law and known
rarity. Independently of that particular sequence,
\leanref{CriticalPsi/Full.lean}{critical\_matching\_psi\_refuted} proves that
for every $C>0$ there are $0<\rho\le1/4$ and $B\ge1$ with
$C\Psi_{\mu,\rho}(B^{-1/2})<R_{\mu,\rho}(B)$.
Its formal witness uses $\rho=\exp(-\exp M)$ and $B=LH/\rho$;
the preceding $B=\rho^{-2}$ sequence remains a paper derivation.
The divergence is slow: $\sqrt H<2.54$ for $\rho\ge10^{-100}$ in the
stated rarity range. Together with the loose proof constants, this
precludes inferring a practically visible separation from the asymptotic
claim alone.
The variational optimizer remains valid. The sharp replacement for
arbitrary laws, Theorem~\ref{thm:entropy-tail-law}, accounts for the
additional statistical difficulty across scales through cost entropy.

\subsection{The formal critical lower bound}
The formal lower proof starts with the following statistical finite-gap result.
Let $0<d_j\le1/128$, $d_j/d_{j+1}\ge2$, and $\log K\ge32$.
One measurable test with an independent seed and observable stopping
time $T$ is run on iid Bernoulli flips of mean $1/2\pm d_j$.
If it has finite expected work and error at most $1/10$ under both
signs for every gap, then
\[
 \sum_{j=1}^K d_j^2\E_{-d_j}T\ge \frac{K\log K}{512}.
\]
\leanref{CriticalEntropy/FiniteGap.lean}{finite\_gap\_sign\_entropy}
derives this from the actual infinite Bernoulli experiment. It assumes
neither a KL bound nor a packing estimate nor termination under the fair
coin. This proves the finite-gap ingredient in
\cref{crit:lem:finite-sign-entropy}.

\leanref{CriticalEntropy/BaseCost.lean}{critical\_base\_cost} now proves
the single-coin bound in \cref{crit:thm:base-cost}: for a universal $c>0$,
\[
 V\E_0T\ge c\rho L(1+\log L),\qquad 0<\rho\le1/4.
\]
The same observable procedure and independent seed law serve every
measurable prior $q\le1/2$ with $\int q\,d\mu\le\rho$. Its uniform
bias is at most $\rho/64$, its extended variance is bounded by $V$,
and expected work is finite under each such prior. The proof derives
square integrability, positive activation, the conditioned procedure,
good gaps, iid restarts, and a sufficient number of scales. A raw
second-moment version is also proved. Neither conditional moment
bounds nor an information inequality is assumed.

The proof then constructs the compact population set, realizes finite
policy mixtures, and obtains one separated pair valid for every capped
policy. Exact finite-prefix simulation laws and work charging transfer
the pair to arbitrary population seeds, coin creation, revisiting and
interleaving. Finite expected work under each endpoint yields termination
and the terminal affinity bound without a continuity assumption on the
output. Thus
\leanref{CriticalEntropy/CriticalMinimax.lean}{critical\_minimax\_lower\_bound}
proves the displayed lower rate for the full $[0,1]$-valued rare-prior
class, including saturation. Separate zero-work results prove that the
budget class is nonempty and its minimax risk is at most $\rho$.
The matching upper and explicit $\Psi$ comparison have the separate formal
proofs cited above; their constructions are described in \cref{sec:verification}.
